% !PRINTED 369
\weberpage{369}
%<<CONT
wir fünf Complexbildungen vom Typus (3). Zwei solche Complexbildungen können sich aber nur dadurch unterscheiden, dass an Stelle von $x_7$ jedesmal ein anderes der Elemente $x_3, x_4, x_5, x_6, x_7$ tritt, und alle diese Möglichkeiten müssen auch vorkommen, weil wir sonst zwei verschiedene Complexe erhalten würden, die dieselben sechs Elemente des vollständigen Systems der $x$ enthalten, was nicht möglich ist, wie wir bewiesen haben. Demnach können wir annehmen, dass die in dem Complexe (3) $\beta$) vorkommenden $x_1, x_2, x_3, x_4, x_5, x_6$ \gesperrt{irgend welche sechs Elemente} des vollständigen Systemes seien, was bewiesen werden sollte.

\S.\,99.

\gesperrt{Der Hesse-Cayley'sche Algorithmus zur Bezeichnung der Doppeltangenten.}

Der Satz 7. des vorigen Paragraphen führt zu einer Bezeichnungsweise für die Doppeltangenten, durch die eine sehr übersichtliche Darstellung aller dieser Verhältnisse möglich ist, die von \gesperrt{Cayley} (im Anschluss an \gesperrt{Hesse}) ausgebildet ist.

Wir legen ein vollständiges Siebener-System

\[ x_1,\ x_2,\ x_3,\ x_4,\ x_5,\ x_6,\ x_7 \]

zu Grunde, dessen Elemente wir einfach durch die Ziffern 1, 2, 3, 4, 5, 6, 7 bezeichnen. Wir sondern eine beliebige Doppeltangente, etwa $x_1$, dieses Systemes aus, und bilden den Complex, der die übrigen sechs enthält:

\weq{1}{x_2 y_2,\ x_3 y_3,\ x_4 y_4,\ x_5 y_5,\ x_6 y_6,\ x_7 y_7.}

Die Doppeltangente $y_2$ ist dann durch das gewählte $x_1$ und durch $x_2$ völlig bestimmt und kann daher durch $[1\,2]$ bezeichnet werden. Ebenso bezeichnen wir $y_3$ durch $[1\,3]$ u.\,s.\,f. Die Doppeltangenten $[1\,2], [1\,3], \dots, [1\,7]$ sind hierdurch vollständig bestimmt und von einander verschieden. Aus dieser Bestimmung geht auch hervor, was man allgemein unter $[\mu\,\nu]$ zu verstehen hat, wenn $\mu, \nu$ zwei verschiedene Ziffern aus der Reihe 1 bis 7 bedeuten.

Es ist nun zunächst zu zeigen, dass $[\mu\,\nu] = [\nu\,\mu]$ ist. Es genügt, wenn wir nachweisen, dass $[1\,2] = [2\,1]$ ist. Dazu brauchen wir nur den Complex zu bilden, der $x_1, x_3, x_4, x_5, x_6, x_7$
%CONT>>
